% ============================================
% Assistant Professor Cover Letter – University of Rhode Island
% ============================================

\documentclass[11pt,letterpaper]{letter}

\usepackage{geometry}
\usepackage{microtype}
\usepackage[hidelinks]{hyperref}

\geometry{left=25mm,right=25mm,top=25mm,bottom=25mm}

\signature{Decory Edwards}
\address{
Department of Economics \\
Johns Hopkins University \\
Baltimore, MD 21218 \\
\texttt{dedwar65@jh.edu}
}

\begin{document}

\begin{letter}{
Search Committee \\
Department of Economics \\
University of Rhode Island \\
Kingston, RI
}

\opening{Dear Members of the Search Committee,}

I am writing to express my strong interest in the tenure-track Assistant Professor position in the Department of Economics at the University of Rhode Island. I am a Ph.D.\ candidate in Economics at Johns Hopkins University, advised by Professors Christopher D.~Carroll, Nicholas W.~Papageorge, and Stelios Fourakis, and I expect to receive my degree in May~2026. My research fields are macroeconomics, wealth and income dynamics, and computational economics.

My research agenda focuses on understanding the determinants of wealth inequality and the mechanisms through which heterogeneity in returns to wealth shapes aggregate outcomes. In my job-market paper, I estimate the degree of heterogeneity in individual portfolio returns using micro-data from the Survey of Consumer Finances and simulate a structural model in which households face idiosyncratic return risk. I show that allowing for persistent heterogeneity in returns substantially improves the model’s ability to match the empirical distribution of wealth, offering a tractable way to connect micro-level return dispersion to macroeconomic inequality.

In a second project, I use the Health and Retirement Study to examine the relationship between interpersonal trust and portfolio performance. Building on the literature that finds a hump-shaped relationship between trust and income, I show that a similar relationship holds for individual returns to wealth measured in the HRS data, highlighting a behavioral channel through which social attitudes shape saving and investment outcomes. This work also provides new estimates of the persistent component of returns to net worth, which motivate the calibration of heterogeneity in my structural model.

The University of Rhode Island’s emphasis on undergraduate and graduate education in a public research university setting strongly appeals to me. I am particularly interested in contributing to a department that serves a diverse student body and values rigorous, policy-relevant economic analysis. My teaching experience includes Principles of Macroeconomics, Intermediate Macroeconomic Theory, and applied macroeconomic topics. In the classroom, I emphasize economic intuition, disciplined use of data, and clear connections between theory and real-world economic issues. I would be well prepared to teach undergraduate and graduate courses aligned with the department’s curricular needs and to mentor students interested in quantitative and policy-oriented research.

Beyond academic fit, I am drawn to URI’s mission as a public land-grant institution with strong ties to the region and a commitment to student access and success. I would be enthusiastic about contributing to the department’s teaching, research, and service missions and about building a long-term academic career at the University of Rhode Island.

Thank you very much for your time and consideration. I would welcome the opportunity to discuss my application further.

\closing{Sincerely,}

\end{letter}
\end{document}
